\section{Структура данных и требования к датасету}

TODO2

\subsection{Форматы входных данных и их препроцессинг}

В проекте ранее уже были практически проверены несколько представлений данных, такие как облака точек XYZ, возникающие при семплировании поверхности или объема, и промежуточный текстовый язык Steps.txt или .json, используемый для описания операций между соседними шагами построения. Но в текущем рабочем подходе используются несколько представлений, дополняющих друг друга. Первое — треугольные сетки STL, на которых можно строить CNN-подобные, графовые или меш-ориентированные модели. Второе — STEP/B-Rep-представление, потенциально пригодное для более глубокого анализа поверхностей, граней, топологии сопряжения и инженерной семантики, и необходимое для автоматической разметки датасета. Третье — промежуточный текстовый структурированный язык в формате JSON, используемый для описания операций между соседними шагами построения.

Таким образом, текущий рабочий подход предполагает использование нескольких представлений в зависимости от стадии конвейера. Для точного восстановления параметров и привязки к инженерному смыслу предпочтительнее STEP/B-Rep. Для разметки датасета и определения пригодности модели для обучения данной операции сохраняется промежуточное командное описание в виде JSON. Для первичного распознавания допустимы STL и производные от него представления.



\subsection{Формирование и очистка обучающей выборки}

Статистика подготовки нового массива данных показывает масштаб инженерной фильтрации. Из 82 418 исходно подготовленных моделей только 44 181 прошли первичную обработку. Далее были выявлены 18 моделей с двумя JSON-файлами, 262 без JSON, 7 417 без STL/STEP и 15 020 с единственной парой STL+STEP, непригодной для обучения переходов между шагами. В итоге рабочий набор составил 21 464 моделей. Эта статистика показывает, что управление качеством данных в проекте является не вспомогательной, а центральной функцией.

На частном эксперименте по операции fillet из всего итогового набора были отобраны только 1 726 моделей, одновременно содержащих скругление и не превышающих порог сложности в 1 500 треугольников. Данный факт важен по двум причинам. Во-первых, большой суммарный датасет не гарантирует достаточный объем примеров для каждой операции. Во-вторых, ограничения по сложности геометрии дополнительно уменьшают фактическую выборку и могут приводить к смещению модели в сторону упрощенных случаев. Следовательно, стратегией развития должен быть не просто рост общего количества моделей, а целевой набор сбалансированных подмассивов по операциям и уровням сложности.

\newpage
\begin{table}[h]
\centering
\caption{Статистика подготовки данных}
\label{tab:data-preparation-stats}
\begin{tabular}{|l|r|p{6cm}|}
\hline
\textbf{Показатель} & \textbf{Значение} & \textbf{Смысл для проекта} \\
\hline
Подготовлено моделей & 82 418 & Сырой массив, отражающий потенциал масштабирования \\
\hline
Прошли первичную обработку & 44 181 & Модели, достигшие стадии первичной машинной обработки \\
\hline
Два JSON & 18 & Структурная неоднозначность метаданных \\
\hline
Нет JSON & 262 & Невозможность supervised-разметки по операциям \\
\hline
Нет STL/STEP & 7 417 & Потеря геометрической основы для обучения \\
\hline
Только одна пара STL+STEP & 15 020 & Непригодность для обучения переходов и шагов \\
\hline
Итоговый рабочий набор & 21 464 & Модели с полным составом данных \\
\hline
Подвыборка для fillet & 1 726 & Фактический обучающий массив для частного сегментационного эксперимента \\
\hline
\end{tabular}
\end{table}
